%===============================================================
% LAPORAN KEMAJUAN MINGGUAN — MINGGU KE-3
% Workshop Pemrograman Framework
% GovSmart — Sistem Pengaduan & Layanan Digital Kabupaten Lamongan
%===============================================================

\documentclass[12pt,a4paper]{article}

% ---- Packages ----
\usepackage[T1]{fontenc}
\usepackage[utf8]{inputenc}
\usepackage[bahasa]{babel}
\usepackage[left=3cm, right=2.5cm, top=2.5cm, bottom=2.5cm]{geometry}
\usepackage{graphicx}
\usepackage{hyperref}
\usepackage{xcolor}
\usepackage{listings}
\usepackage{enumitem}
\usepackage{booktabs}
\usepackage{longtable}
\usepackage{array}
\usepackage{fancyhdr}
\usepackage{titlesec}
\usepackage{tcolorbox}
\usepackage{fontawesome5}
\usepackage{parskip}
\usepackage{microtype}
\usepackage{caption}
\usepackage{float}
\usepackage{multirow}

% ---- Warna ----
\definecolor{accent}{HTML}{1E40AF}
\definecolor{accentlight}{HTML}{EFF6FF}
\definecolor{success}{HTML}{059669}
\definecolor{successlight}{HTML}{ECFDF5}
\definecolor{warning}{HTML}{D97706}
\definecolor{warninglight}{HTML}{FFFBEB}
\definecolor{codebg}{HTML}{F8FAFC}
\definecolor{codefg}{HTML}{1E293B}
\definecolor{borderblue}{HTML}{BFDBFE}
\definecolor{gray}{HTML}{64748B}

% ---- Hyperref ----
\hypersetup{
  colorlinks=true,
  linkcolor=accent,
  urlcolor=accent,
  pdftitle={Laporan Kemajuan Minggu 3 - GovSmart},
  pdfauthor={Moh. Faiz Khoiru Izzuddin},
  pdfsubject={Workshop Pemrograman Framework},
}

% ---- Header / Footer ----
\pagestyle{fancy}
\fancyhf{}
\fancyhead[L]{\small\textcolor{gray}{GovSmart --- Laporan Kemajuan Minggu 3}}
\fancyhead[R]{\small\textcolor{gray}{Workshop Pemrograman Framework}}
\fancyfoot[C]{\small\textcolor{gray}{\thepage}}
\renewcommand{\headrulewidth}{0.5pt}
\renewcommand{\headrule}{\hbox to\headwidth{\color{borderblue}\leaders\hrule height \headrulewidth\hfill}}

% ---- Section Format ----
\titleformat{\section}
  {\large\bfseries\color{accent}}
  {\thesection.}{0.75em}{}
  [\vspace{2pt}\color{borderblue}\hrule\vspace{4pt}]

\titleformat{\subsection}
  {\normalsize\bfseries\color{accent!80!black}}
  {\thesubsection.}{0.6em}{}

\titleformat{\subsubsection}
  {\normalsize\bfseries\color{codefg}}
  {\thesubsubsection.}{0.5em}{}

% ---- Listing / Code ----
\lstset{
  backgroundcolor=\color{codebg},
  basicstyle=\ttfamily\footnotesize\color{codefg},
  breaklines=true,
  breakatwhitespace=true,
  frame=single,
  framesep=6pt,
  rulecolor=\color{borderblue},
  xleftmargin=12pt,
  xrightmargin=6pt,
  aboveskip=10pt,
  belowskip=6pt,
  showstringspaces=false,
  tabsize=2,
  keywordstyle=\color{accent}\bfseries,
  commentstyle=\color{gray}\itshape,
  stringstyle=\color{success},
  numbers=left,
  numberstyle=\tiny\color{gray},
  numbersep=8pt,
}

% ---- tcolorbox styles ----
\tcbuselibrary{skins,breakable}

\newtcolorbox{infobox}[1][]{
  enhanced,
  colback=accentlight,
  colframe=accent,
  fonttitle=\bfseries\small,
  title=#1,
  arc=6pt,
  left=8pt, right=8pt, top=6pt, bottom=6pt,
  breakable
}

\newtcolorbox{successbox}[1][]{
  enhanced,
  colback=successlight,
  colframe=success,
  fonttitle=\bfseries\small,
  title=#1,
  arc=6pt,
  left=8pt, right=8pt, top=6pt, bottom=6pt,
  breakable
}

\newtcolorbox{warningbox}[1][]{
  enhanced,
  colback=warninglight,
  colframe=warning,
  fonttitle=\bfseries\small,
  title=#1,
  arc=6pt,
  left=8pt, right=8pt, top=6pt, bottom=6pt,
  breakable
}

% ---- Dokumen ----
\begin{document}

%------------------------------------------------------------
% HALAMAN JUDUL
%------------------------------------------------------------
\begin{titlepage}
  \centering
  \vspace*{1.5cm}

  {\Huge\bfseries\color{accent} LAPORAN KEMAJUAN\\[6pt]
   \Large WORKSHOP PEMROGRAMAN FRAMEWORK}\\[0.4cm]
  {\large\color{gray}\rule{12cm}{1.5pt}}\\[0.6cm]

  {\LARGE\bfseries GovSmart}\\[0.2cm]
  {\large\color{gray}
    Sistem Pengaduan \& Layanan Digital\\
    Pemerintah Kabupaten Lamongan}\\[0.8cm]

  \begin{tcolorbox}[
    enhanced, width=10cm,
    colback=accentlight, colframe=accent,
    arc=10pt, boxrule=1.2pt,
    fontupper=\bfseries\large\color{accent},
    halign=center
  ]
    PERKEMBANGAN MINGGU KE-3
  \end{tcolorbox}

  \vspace{1.5cm}

  \begin{tabular}{ll}
    \textbf{Nama}          & : Moh.\ Faiz Khoiru Izzuddin \\[4pt]
    \textbf{NIM}           & : 3125521049 \\[4pt]
    \textbf{Program Studi} & : Teknik Informatika \\[4pt]
    \textbf{Mata Kuliah}   & : Workshop Pemrograman Framework \\[4pt]
    \textbf{Periode}       & : Minggu Ke-3 \\[4pt]
    \textbf{Tanggal}       & : Oktober 2026 \\
  \end{tabular}

  \vfill
  {\small\color{gray}
    Laporan ini merupakan dokumentasi perkembangan proyek GovSmart\\
    sebagai bagian dari tugas Workshop Pemrograman Framework.}
\end{titlepage}

%------------------------------------------------------------
% DAFTAR ISI
%------------------------------------------------------------
\tableofcontents
\newpage

%------------------------------------------------------------
% 1. RINGKASAN EKSEKUTIF
%------------------------------------------------------------
\section{Ringkasan Eksekutif}

Pada minggu ketiga pengerjaan proyek \textbf{GovSmart}, fokus utama diarahkan pada
\textbf{penyelesaian menyeluruh lapisan antarmuka pengguna (\textit{frontend})} berbasis
React + Vite, setelah fondasi backend FastAPI berhasil dibangun pada minggu-minggu sebelumnya.
Seluruh fitur interaktif yang direncanakan kini telah terimplementasi secara penuh dan
terintegrasi dengan API backend melalui mekanisme \textit{reverse proxy} Vite.

\begin{successbox}[Pencapaian Utama Minggu 3]
  \begin{itemize}[leftmargin=1.5em, itemsep=3pt]
    \item \textbf{Portal Autentikasi Glassmorphism} --- Halaman login/registrasi
          dengan animasi sliding panel selesai diimplementasikan.
    \item \textbf{Halaman Beranda Hero MSN-Style} --- Tampilan utama bergaya portal
          berita modern dengan latar gambar Lamongan.
    \item \textbf{Modul Pengaduan Warga Lengkap} --- Form pengaduan dengan GPS
          auto-fill, drag-and-drop upload foto, dan feed laporan publik.
    \item \textbf{Integrasi Peta Interaktif Leaflet} --- Peta sebaran laporan
          dengan pin warna per dinas.
    \item \textbf{Dasbor Admin \& Super Admin} --- Manajemen laporan, ubah status,
          audit log, CRUD admin dinas.
    \item \textbf{AI Chatbot Widget} --- \textit{Floating chatbot} terhubung ke
          endpoint AI backend.
    \item \textbf{Sistem Notifikasi Toast} --- Umpan balik real-time untuk setiap
          aksi pengguna.
    \item \textbf{Desain Responsif} --- Media queries untuk tampilan mobile dan tablet.
  \end{itemize}
\end{successbox}

\vspace{0.5em}
Total baris kode yang ditulis pada minggu ini mencapai lebih dari \textbf{2.900 baris}
(1.908 baris JSX di \texttt{App.jsx} + 1.000 baris CSS di \texttt{index.css}), menjadikan
minggu ini sebagai periode paling produktif dalam siklus pengembangan frontend proyek ini.

%------------------------------------------------------------
% 2. RECAP MINGGU SEBELUMNYA
%------------------------------------------------------------
\section{Rekap Perkembangan Minggu Sebelumnya}

Sebelum membahas pencapaian minggu ketiga, berikut adalah ringkasan kondisi proyek
yang telah dicapai hingga akhir minggu kedua:

\begin{infobox}[Status Akhir Minggu 2]
  \begin{itemize}[leftmargin=1.5em, itemsep=3pt]
    \item Backend FastAPI telah berjalan lengkap dengan endpoint autentikasi JWT,
          laporan, dinas, admin, chatbot AI, dan audit log.
    \item Skema database PostgreSQL telah selesai dengan tabel \texttt{users},
          \texttt{reports}, \texttt{departments}, dan \texttt{audit\_logs}.
    \item Integrasi AI (Gemini API) untuk \textit{smart routing} laporan dan chatbot
          sudah aktif.
    \item Frontend masih dalam kondisi awal/scaffolding --- belum memiliki tampilan
          yang lengkap dan interaktif.
    \item Kontainerisasi Docker sudah dikonfigurasi namun belum diuji penuh.
  \end{itemize}
\end{infobox}

%------------------------------------------------------------
% 3. IMPLEMENTASI FRONTEND MINGGU 3
%------------------------------------------------------------
\section{Implementasi Frontend Minggu Ke-3}

\subsection{Arsitektur dan Struktur File Frontend}

Frontend GovSmart menggunakan pendekatan \textit{single-file component} dimana seluruh
logika, state management, dan tampilan dikelola dalam satu file utama \texttt{App.jsx},
didukung oleh \texttt{index.css} untuk semua style global.

\begin{lstlisting}[language=bash, caption=Struktur direktori frontend]
frontend/
  src/
    App.jsx        # Komponen utama (1908 baris JSX)
    index.css      # Global CSS & Design System (1000 baris)
    main.jsx       # Entry point React
  public/
    bg-lamongan.jpg  # Latar gambar hero
  vite.config.js     # Konfigurasi Vite + Proxy API
  package.json
\end{lstlisting}

\subsubsection{Dependensi Utama Frontend}

\begin{center}
\begin{tabular}{llp{6.5cm}}
  \toprule
  \textbf{Paket} & \textbf{Versi} & \textbf{Fungsi} \\
  \midrule
  react & 18.x & Library UI utama \\
  react-dom & 18.x & Rendering DOM \\
  vite & 5.x & Build tool \& dev server \\
  lucide-react & latest & Icon SVG modern \\
  react-leaflet & 4.x & Komponen peta interaktif \\
  leaflet & 1.9.4 & Peta OpenStreetMap \\
  \bottomrule
\end{tabular}
\end{center}

\subsection{Design System dan CSS Global}

Seluruh desain dibangun di atas \textbf{Design System berbasis CSS Variables} yang
didefinisikan dalam \texttt{:root} di \texttt{index.css}. Pendekatan ini memungkinkan
konsistensi warna, ukuran, dan bayangan di seluruh komponen.

\begin{lstlisting}[language=CSS, caption=CSS Variables (Design Tokens) pada index.css]
:root {
  --bg-primary:       #f0f4f8;
  --bg-surface:       #ffffff;
  --accent-primary:   #1e40af;   /* Biru pemerintah */
  --accent-secondary: #b45309;
  --text-main:        #000000;
  --text-muted:       #1f2937;
  --shadow-glow:      0 4px 12px rgba(30, 64, 175, 0.2);
  --shadow-lg:        0 10px 25px -5px rgba(0,0,0,0.15);
  --radius-xl:        16px;
  --radius-lg:        12px;
}
\end{lstlisting}

Tipografi menggunakan \textbf{Outfit} dari Google Fonts yang diimpor secara dinamis,
memberikan tampilan modern dan profesional yang sesuai untuk portal pemerintahan.

\begin{lstlisting}[language=CSS, caption=Import Font Outfit]
@import url("https://fonts.googleapis.com/css2?
  family=Outfit:wght@300;400;500;600;700;800&display=swap");

* {
  box-sizing: border-box;
  font-family: "Outfit", sans-serif;
}
\end{lstlisting}

Latar belakang menggunakan teknik \textbf{Fixed Background Image with Gradient Overlay},
menampilkan foto Kabupaten Lamongan yang tetap meski halaman digulir:

\begin{lstlisting}[language=CSS, caption=Latar belakang halaman utama]
body {
  background-image:
    linear-gradient(
      to bottom,
      rgba(0,0,0,0.6) 0%,
      rgba(0,0,0,0) 30%,
      rgba(0,0,0,0.1) 100%
    ),
    url("/bg-lamongan.jpg");
  background-size: cover;
  background-position: center;
  background-attachment: fixed;
}
\end{lstlisting}

\subsection{Modul Autentikasi (Login \& Registrasi)}

Portal autentikasi diimplementasikan dengan desain \textbf{Glassmorphism} dua panel
yang menampilkan panel visual kiri (branding GovSmart) dan panel kanan (form interaktif).
Peralihan antara mode \textit{login} dan \textit{register} menggunakan animasi
\textbf{sliding CSS} tanpa library tambahan.

\subsubsection{Fitur Portal Autentikasi}

\begin{itemize}[leftmargin=2em, itemsep=4pt]
  \item \textbf{Dua Panel Split:} Panel kiri berisi branding dan statistik layanan;
        panel kanan berisi form autentikasi.
  \item \textbf{Animasi Sliding:} Transisi halus antar form login dan register
        menggunakan \texttt{translateX} CSS dan state \texttt{authMode}.
  \item \textbf{Validasi Error:} Pesan error detail dari backend ditampilkan secara
        user-friendly (mis.\ validasi NIK, email, password minimal karakter).
  \item \textbf{Auto-Login Pasca Register:} Setelah registrasi sukses, pengguna
        diarahkan ke mode login secara otomatis.
  \item \textbf{Persistent Session:} Token JWT disimpan di \texttt{localStorage}
        dan di-\textit{restore} saat halaman dimuat ulang.
  \item \textbf{Ambient Glow Effect:} Dua elemen dekoratif berbentuk lingkaran
        blur berwarna kuning dan hijau sebagai efek ambient.
\end{itemize}

\begin{lstlisting}[language=CSS, caption=Sliding panel transition CSS]
.auth-form-panel {
  position: absolute;
  inset: 0;
  opacity: 0;
  transform: translateX(32px);
  transition:
    transform 0.48s cubic-bezier(0.22, 1, 0.36, 1),
    opacity 0.34s ease,
    visibility 0.34s ease;
}

/* Panel aktif tampil, panel non-aktif geser keluar */
.auth-form-track:not(.is-register) .auth-form-panel:first-child,
.auth-form-track.is-register .auth-form-panel:last-child {
  opacity: 1;
  transform: translateX(0);
}
\end{lstlisting}

\subsubsection{Alur Autentikasi Frontend}

\begin{lstlisting}[language=JavaScript, caption=Handler login dengan JWT dan role detection]
const handleLogin = async (e) => {
  e.preventDefault();
  // 1. Kirim credentials ke endpoint OAuth2
  const res = await fetch("/api/v1/auth/login", {
    method: "POST",
    headers: { "Content-Type": "application/x-www-form-urlencoded" },
    body: formData,
  });
  const data = await res.json();

  // 2. Simpan token ke localStorage
  localStorage.setItem("govsmart_token", data.access_token);
  setToken(data.access_token);

  // 3. Deteksi role melalui endpoint /me
  const profileRes = await fetch("/api/v1/admin/me", {
    headers: { Authorization: `Bearer ${data.access_token}` },
  });
  const profile = await profileRes.json();
  setUserRole(profile.role); // 'warga' | 'admin' | 'super_admin'
  setUserEmail(profile.email);
};
\end{lstlisting}

\subsection{Halaman Beranda Hero (MSN-Style)}

Setelah login berhasil, pengguna disambut dengan tampilan \textbf{Hero Section}
bergaya portal berita modern ala MSN. Halaman ini hanya muncul saat
\texttt{mainView === "options"}.

\subsubsection{Elemen Halaman Beranda}

\begin{enumerate}[leftmargin=2em, itemsep=4pt]
  \item \textbf{Navbar Solid Putih:} Berisi logo GovSmart AI, info pengguna yang
        sedang login (email + role), dan tombol Keluar.
  \item \textbf{Hero Section Transparan:} Judul besar dengan \textit{text-shadow}
        tebal agar terbaca di atas gambar Lamongan.
  \item \textbf{Search Bar Dekoratif:} \textit{Pill-shaped} search bar putih
        dengan bayangan besar sebagai elemen visual.
  \item \textbf{Quick Links Row:} Dua atau satu ikon navigasi cepat berbentuk
        lingkaran dengan latar \textit{frosted glass} --- navigasi ke Pengaduan
        atau Pertanyaan (untuk warga), atau ke Dasbor Admin (untuk admin/super\_admin).
  \item \textbf{Content Board Melengkung:} Panel putih besar melengkung di bagian
        bawah sebagai transisi dari hero ke konten.
\end{enumerate}

\begin{lstlisting}[language=CSS, caption=Animasi entrance fade-in-up]
@keyframes fade-in-up {
  from { opacity: 0; transform: translateY(20px); }
  to   { opacity: 1; transform: translateY(0); }
}

.hero-msn    { animation: fade-in-up 0.6s ease-out; }
.content-board { animation: fade-in-up 0.8s ease-out; }
\end{lstlisting}

\subsection{Modul Pengaduan Warga}

Modul pengaduan adalah fitur inti aplikasi yang dapat diakses oleh warga terdaftar.
Halaman ini menggunakan layout \textbf{dua kolom} (grid 1fr 1fr): kolom kiri berisi
form pengaduan, kolom kanan menampilkan feed laporan publik.

\subsubsection{Form Pengaduan Baru}

Form pengaduan memiliki fitur-fitur berikut:

\begin{enumerate}[leftmargin=2em, itemsep=6pt]

  \item \textbf{Input Judul \& Deskripsi:} Field teks standar dengan validasi
        HTML5 (\texttt{required}).

  \item \textbf{GPS Auto-Fill Koordinat:}
        Tombol \textit{``Gunakan Lokasi Saya''} memanfaatkan Geolocation API browser
        untuk mengisi otomatis latitude dan longitude dengan presisi 6 desimal.
        \begin{lstlisting}[language=JavaScript, caption=Handler GPS auto-fill]
const handleUseGPS = () => {
  setGpsLoading(true);
  navigator.geolocation.getCurrentPosition(
    (pos) => {
      setLat(pos.coords.latitude.toFixed(6));
      setLong(pos.coords.longitude.toFixed(6));
      setGpsLoading(false);
    },
    () => {
      notify("Gagal mendapatkan lokasi.");
      setGpsLoading(false);
    }
  );
};
        \end{lstlisting}

  \item \textbf{Drag \& Drop Upload Foto:}
        Area unggah foto mendukung \textit{drag-and-drop} dengan visual feedback
        (border dan background berubah saat file di-drag). File diunggah ke
        endpoint \texttt{/api/v1/reports/\{id\}/upload-foto} secara otomatis
        setelah laporan berhasil dibuat.
        \begin{lstlisting}[language=JavaScript, caption=Drag-drop handler]
onDrop={(e) => {
  e.preventDefault();
  setIsDragging(false);
  const file = e.dataTransfer.files[0];
  if (file && ["image/jpeg","image/png","image/webp"]
      .includes(file.type)) {
    setFotoFile(file);
  }
}}
        \end{lstlisting}

  \item \textbf{Konfirmasi Sukses dengan Info AI:}
        Setelah laporan dikirim, ditampilkan panel konfirmasi hijau berisi
        tingkat urgensi yang dianalisis AI (darurat/penting/normal) dan
        nama dinas yang menerima laporan (hasil \textit{smart routing} AI).

\end{enumerate}

\subsubsection{Feed Pengaduan Publik}

Panel kanan menampilkan daftar semua laporan yang ada di sistem secara \textit{real-time}
(di-refresh otomatis setelah submit). Setiap kartu laporan menampilkan:
badge urgensi berwarna, judul laporan, isi ringkas, nama dinas tujuan, status terkini,
dan tanggal dibuat.

\begin{lstlisting}[language=JavaScript, caption=Fetch laporan dari API]
const fetchReports = async () => {
  setLoadingReports(true);
  const headers = token
    ? { Authorization: `Bearer ${token}` }
    : {};
  const res = await fetch("/api/v1/reports/", { headers });
  if (res.ok) setReports(await res.json());
  setLoadingReports(false);
};
\end{lstlisting}

\subsubsection{Badge Status Urgensi}

\begin{center}
\begin{tabular}{lll}
  \toprule
  \textbf{Nilai Urgensi} & \textbf{Warna Latar} & \textbf{Warna Teks} \\
  \midrule
  darurat & Merah muda (\#fff1f2) & Merah (\#e11d48) \\
  penting & Kuning muda (\#fffbeb) & Amber (\#d97706) \\
  normal  & Hijau muda (\#ecfdf5) & Hijau (\#059669) \\
  \bottomrule
\end{tabular}
\end{center}

\subsection{Modul Pertanyaan --- AI Chatbot Halaman Penuh}

Modul Pertanyaan menampilkan antarmuka chatbot AI dalam mode halaman penuh
(bukan widget \textit{floating}). Fitur ini dapat diakses warga untuk bertanya
mengenai panduan layanan publik, prosedur perizinan UMKM/NIB, dan informasi dinas.

\begin{itemize}[leftmargin=2em, itemsep=4pt]
  \item Area pesan bergulir dengan tinggi tetap 400px.
  \item Bubble chat user (kanan, biru) dan AI (kiri, putih).
  \item \textbf{Typing Indicator:} Tiga titik animasi bounce saat AI sedang memproses.
  \item Pesan awal dari AI muncul otomatis saat halaman dibuka.
\end{itemize}

\begin{lstlisting}[language=CSS, caption=Animasi typing indicator]
.typing-dot {
  width: 8px; height: 8px;
  background: var(--text-dim);
  border-radius: 50%;
  animation: typing-bounce 1.4s infinite ease-in-out both;
}
@keyframes typing-bounce {
  0%, 80%, 100% { transform: scale(0); }
  40%           { transform: scale(1); }
}
\end{lstlisting}

\subsection{Dasbor Admin \& Super Admin}

Pengguna dengan role \texttt{admin} atau \texttt{super\_admin} diarahkan ke
\textbf{Dasbor Manajemen Pengaduan} yang terdiri dari beberapa panel:

\subsubsection{Panel 1: Registrasi Admin Dinas (Super Admin only)}

Hanya dapat diakses oleh \texttt{super\_admin}. Form 4-kolom untuk membuat akun
admin baru: NIP Aparatur, Email Dinas, Password, dan pilihan Dinas. Data dikirim
ke endpoint \texttt{POST /api/v1/admin/admin-dinas}.

\subsubsection{Panel 2: Daftar Admin Dinas Terdaftar (Super Admin only)}

Tabel daftar semua admin dinas yang terdaftar dengan kolom ID, NIP, Email, Dinas,
Status (Aktif/Non-aktif), dan Aksi (Edit/Hapus).

\begin{itemize}[leftmargin=2em, itemsep=4pt]
  \item \textbf{Edit Admin:} Modal popup untuk mengubah NIP, email, password, dan dinas.
  \item \textbf{Hapus/Nonaktifkan Admin:} Konfirmasi \texttt{window.confirm()} sebelum
        pemanggilan \texttt{DELETE /api/v1/admin/admin-dinas/\{id\}}.
\end{itemize}

\subsubsection{Panel 3: Peta Sebaran Laporan Interaktif}

Peta interaktif menggunakan \textbf{React Leaflet} dengan tile dari OpenStreetMap,
berpusat pada koordinat Kabupaten Lamongan \texttt{[-7.1189, 112.4109]}.

\begin{itemize}[leftmargin=2em, itemsep=4pt]
  \item Setiap laporan yang memiliki koordinat ditampilkan sebagai pin khusus.
  \item Pin berwarna berbeda per dinas (8 warna dari palet \texttt{DEPT\_COLORS}).
  \item \textbf{Custom Marker:} SVG ikon dokumen berbentuk teardrop dengan warna
        dinas masing-masing.
  \item \textbf{Popup Info:} Klik pin menampilkan judul laporan, nama dinas,
        tingkat urgensi, dan status.
  \item \textbf{Legenda Peta:} Deretan badge warna-dinas di atas peta.
\end{itemize}

\begin{lstlisting}[language=JavaScript, caption=Custom Leaflet marker per dinas]
const DEPT_COLORS = [
  "#ef4444", "#f97316", "#eab308", "#22c55e",
  "#14b8a6", "#3b82f6", "#8b5cf6", "#ec4899"
];

const createReportIcon = (color) =>
  L.divIcon({
    className: "report-map-marker",
    html: `<div style="--marker-color:${color}">
             <svg viewBox="0 0 24 24">
               <path d="M14 2H6a2 2 0 0 0-2 2v16..." />
             </svg>
           </div>`,
    iconSize: [32, 40],
    iconAnchor: [16, 40],
    popupAnchor: [0, -38],
  });
\end{lstlisting}

\subsubsection{Panel 4: Tabel Daftar Pengaduan Masuk}

Tabel komprehensif semua laporan yang masuk dengan kolom ID, Urgensi AI,
Judul Laporan, Status, dan Aksi. Setiap baris memiliki dua tombol aksi:

\begin{itemize}[leftmargin=2em, itemsep=4pt]
  \item \textbf{Ubah Status:} Membuka modal untuk mengubah status laporan
        (diterima \textrightarrow{} diproses \textrightarrow{} selesai / ditolak).
  \item \textbf{Audit Log:} Membuka modal riwayat perubahan status laporan
        lengkap dengan timestamp dan ID admin yang mengubah.
\end{itemize}

\subsection{Modal Interaktif}

Tiga jenis modal diimplementasikan dengan komponen yang konsisten:

\begin{center}
\begin{longtable}{p{3.5cm}p{4cm}p{5cm}}
  \toprule
  \textbf{Modal} & \textbf{Trigger} & \textbf{Fungsi} \\
  \midrule
  Ubah Status & Tombol ``Ubah Status'' & Dropdown pilihan status baru \\
  Audit Log & Tombol ``Audit Log'' & Riwayat perubahan status + timestamp \\
  Edit Admin Dinas & Tombol ``Edit'' di tabel admin & Form edit NIP, email, password, dinas \\
  \bottomrule
\end{longtable}
\end{center}

\subsection{AI Chatbot Widget (Floating)}

Di samping halaman Pertanyaan, GovSmart juga menyediakan \textbf{floating chatbot widget}
yang dapat diakses dari halaman mana pun setelah login. Widget ini berada di pojok
kanan bawah layar.

\begin{itemize}[leftmargin=2em, itemsep=4pt]
  \item Tombol bulat biru 64px dengan ikon Bot dari Lucide.
  \item Panel chat 400x600px muncul di atas tombol dengan animasi.
  \item Pesan AI dan pengguna ditampilkan dalam gelembung chat.
  \item Terhubung ke endpoint \texttt{POST /api/v1/chat/}.
  \item Token Bearer disertakan secara otomatis jika pengguna sudah login.
\end{itemize}

\begin{lstlisting}[language=CSS, caption=Styling chatbot trigger button]
.chatbot-trigger {
  position: fixed;
  bottom: 2.5rem;
  right: 2.5rem;
  width: 64px;
  height: 64px;
  border-radius: 50%;
  background: var(--accent-primary);
  box-shadow: 0 10px 25px rgba(30, 64, 175, 0.4);
  z-index: 90;
}
.chatbot-trigger:hover {
  transform: scale(1.05);
  background: #1e3a8a;
}
\end{lstlisting}

\subsection{Sistem Notifikasi Toast}

Sistem notifikasi \textit{toast} non-blocking diimplementasikan menggunakan state
React tanpa dependensi library eksternal. Toast muncul di pojok kanan atas dengan
animasi slide-in dan hilang otomatis setelah 4,5 detik.

\begin{itemize}[leftmargin=2em, itemsep=4pt]
  \item \textbf{Toast Sukses:} Border kiri hijau, ikon CheckCircle.
  \item \textbf{Toast Error:} Border kiri merah, ikon AlertTriangle.
  \item Dapat ditutup manual dengan tombol $\times$.
  \item Mendukung beberapa toast sekaligus (stack vertikal).
\end{itemize}

\begin{lstlisting}[language=JavaScript, caption=Fungsi notify dengan auto-dismiss]
const notify = (message, type = "error") => {
  const id = Date.now() + Math.random();
  setNotifications((items) =>
    [...items, { id, message, type }]);
  setTimeout(
    () => setNotifications((items) =>
      items.filter((item) => item.id !== id)),
    4500
  );
};
\end{lstlisting}

\begin{lstlisting}[language=CSS, caption=Animasi toast masuk]
@keyframes toast-in {
  from { opacity: 0; transform: translateY(-8px) scale(0.98); }
  to   { opacity: 1; transform: translateY(0) scale(1); }
}
.toast { animation: toast-in 0.25s ease-out; }
\end{lstlisting}

\subsection{Desain Responsif}

Semua komponen utama telah dioptimalkan untuk berbagai ukuran layar menggunakan
\textit{CSS Media Queries}:

\begin{center}
\begin{tabular}{lll}
  \toprule
  \textbf{Breakpoint} & \textbf{Target} & \textbf{Penyesuaian} \\
  \midrule
  \texttt{max-width: 860px} & Tablet & Portal auth 1-kolom, padding berkurang \\
  \texttt{max-width: 520px} & Mobile & Border-radius hilang, padding minimal \\
  \bottomrule
\end{tabular}
\end{center}

%------------------------------------------------------------
% 4. STATE MANAGEMENT
%------------------------------------------------------------
\section{Manajemen State React}

Seluruh state aplikasi dikelola menggunakan \texttt{useState} dan \texttt{useEffect}
dari React hooks. Tidak ada library state management eksternal (seperti Redux)
yang digunakan --- hal ini menjaga ukuran \textit{bundle} tetap kecil.

\begin{center}
\begin{longtable}{p{4.5cm}p{7cm}}
  \toprule
  \textbf{State Variable} & \textbf{Kegunaan} \\
  \midrule
  \texttt{mainView} & Navigasi antar halaman (options/pengaduan/pertanyaan/admin) \\
  \texttt{token, userRole, userEmail} & Data sesi pengguna yang login \\
  \texttt{reports, departments} & Data laporan dan dinas dari API \\
  \texttt{notifications} & Antrian pesan toast \\
  \texttt{selectedReport, showAuditModal} & State modal admin \\
  \texttt{lat, long, gpsLoading} & State GPS auto-fill \\
  \texttt{fotoFile, isDragging} & State drag-and-drop foto \\
  \texttt{chatMessages, chatLoading} & State percakapan chatbot \\
  \texttt{adminDinasList, editingAdmin} & State manajemen admin dinas \\
  \bottomrule
\end{longtable}
\end{center}

%------------------------------------------------------------
% 5. INTEGRASI BACKEND
%------------------------------------------------------------
\section{Integrasi Frontend--Backend}

Komunikasi antara frontend React dan backend FastAPI dilakukan melalui
\textbf{Fetch API} bawaan browser. Konfigurasi \textit{proxy} di \texttt{vite.config.js}
meneruskan semua request \texttt{/api/*} ke server backend:

\begin{lstlisting}[language=JavaScript, caption=Vite proxy config untuk API]
// vite.config.js
export default defineConfig({
  server: {
    proxy: {
      '/api': {
        target: 'http://localhost:8000',
        changeOrigin: true,
      }
    }
  }
})
\end{lstlisting}

\begin{center}
\begin{longtable}{p{3cm}p{3cm}p{6cm}}
  \toprule
  \textbf{Endpoint} & \textbf{Metode} & \textbf{Digunakan Oleh} \\
  \midrule
  \texttt{/api/v1/auth/login} & POST & handleLogin \\
  \texttt{/api/v1/auth/register} & POST & handleRegister \\
  \texttt{/api/v1/admin/me} & GET & Deteksi role setelah login \\
  \texttt{/api/v1/reports/} & GET / POST & fetchReports, handleSubmitReport \\
  \texttt{/api/v1/reports/\{id\}/upload-foto} & POST & Upload foto bukti \\
  \texttt{/api/v1/reports/\{id\}/status} & PATCH & handleUpdateStatus \\
  \texttt{/api/v1/departments/} & GET & fetchDepartments \\
  \texttt{/api/v1/admin/admin-dinas} & GET / POST & Manajemen admin dinas \\
  \texttt{/api/v1/admin/admin-dinas/\{id\}} & PUT / DELETE & Edit, hapus admin \\
  \texttt{/api/v1/admin/audit-logs/\{id\}} & GET & handleViewAuditLogs \\
  \texttt{/api/v1/chat/} & POST & handleSendChat \\
  \bottomrule
\end{longtable}
\end{center}

%------------------------------------------------------------
% 6. KENDALA & SOLUSI
%------------------------------------------------------------
\section{Kendala yang Dihadapi dan Solusi}

\begin{enumerate}[leftmargin=2em, itemsep=10pt]

  \item \textbf{Leaflet Default Icon Broken pada Bundler Vite}\\
        Ikon marker bawaan Leaflet tidak muncul karena Vite me-resolve URL aset
        secara berbeda dibanding webpack. \\
        \textit{Solusi:} Menghapus \texttt{\_getIconUrl} dari prototype dan
        men-override ikon secara eksplisit menggunakan CDN unpkg:
        \begin{lstlisting}[language=JavaScript]
delete L.Icon.Default.prototype._getIconUrl;
L.Icon.Default.mergeOptions({
  iconUrl: "https://unpkg.com/leaflet@1.9.4/dist/images/marker-icon.png",
  shadowUrl: "https://unpkg.com/leaflet@1.9.4/dist/images/marker-shadow.png",
});
        \end{lstlisting}

  \item \textbf{Autofill Browser Mengubah Warna Input}\\
        Browser Chromium menerapkan warna latar kuning pada field yang di-autofill,
        merusak tampilan form autentikasi. \\
        \textit{Solusi:} Menggunakan pseudo-selector \texttt{:-webkit-autofill}
        dengan override warna paksa:
        \begin{lstlisting}[language=CSS]
.auth-input:-webkit-autofill {
  -webkit-text-fill-color: #000000 !important;
  background-color: #ffffff !important;
}
        \end{lstlisting}

  \item \textbf{Upload Foto Bergantung pada ID Laporan yang Baru Dibuat}\\
        API upload foto membutuhkan ID laporan yang baru dibuat. \\
        \textit{Solusi:} Upload foto dilakukan secara berurutan setelah respons
        \texttt{POST /reports/} sukses, menggunakan \texttt{data.id} dari respons.

  \item \textbf{CORS pada Development}\\
        Request dari \texttt{localhost:5173} ke \texttt{localhost:8000} diblokir CORS. \\
        \textit{Solusi:} Menggunakan fitur proxy Vite sehingga semua request ke
        \texttt{/api} diteruskan oleh Vite server tanpa masalah CORS.

\end{enumerate}

%------------------------------------------------------------
% 7. RINGKASAN TAMPILAN
%------------------------------------------------------------
\section{Ringkasan Tampilan Halaman Aplikasi}

\begin{center}
\begin{longtable}{p{3.8cm}p{3.8cm}p{5cm}}
  \toprule
  \textbf{Halaman/Modul} & \textbf{Role yang Bisa Akses} & \textbf{Elemen Utama} \\
  \midrule
  Login / Register & Semua (belum login) & 2-panel glassmorphism, sliding animation \\
  Beranda Pilihan & Semua yang sudah login & Hero foto Lamongan, quick links \\
  Form Pengaduan & Warga & Form + GPS + Drag-drop foto + Feed publik \\
  Pertanyaan AI & Warga & Chatbot full-page \\
  Dasbor Admin & Admin, Super Admin & Tabel laporan, peta, ubah status \\
  Dasbor Super Admin & Super Admin saja & + Form \& tabel manajemen Admin Dinas \\
  \bottomrule
\end{longtable}
\end{center}

%------------------------------------------------------------
% 8. RENCANA MINGGU BERIKUTNYA
%------------------------------------------------------------
\section{Rencana Pengembangan Minggu Berikutnya}

Dengan selesainya implementasi frontend secara menyeluruh, rencana pada minggu
keempat mencakup:

\begin{enumerate}[leftmargin=2em, itemsep=6pt]
  \item \textbf{Pengujian Integrasi End-to-End:} Pengujian alur lengkap dari
        registrasi warga, pengiriman laporan, hingga penanganan oleh admin.
  \item \textbf{Penyempurnaan UX:} Menambahkan \textit{loading skeleton} pada
        feed laporan dan tabel admin saat data sedang di-fetch.
  \item \textbf{Optimasi Performa:} Lazy loading komponen berat seperti peta
        Leaflet menggunakan \texttt{React.lazy()} dan \texttt{Suspense}.
  \item \textbf{Pengujian Responsivitas:} Pengujian di berbagai perangkat
        (mobile, tablet) dan perbaikan bug tampilan.
  \item \textbf{Deployment \& Docker:} Menguji penuh konfigurasi Docker Compose
        (frontend + backend + PostgreSQL) dan deployment ke server.
  \item \textbf{Dokumentasi Teknis:} Melengkapi dokumentasi API dan panduan
        instalasi proyek.
\end{enumerate}

%------------------------------------------------------------
% 9. KESIMPULAN
%------------------------------------------------------------
\section{Kesimpulan}

Minggu ketiga merupakan fase implementasi paling intensif dalam siklus pengembangan
proyek GovSmart. Seluruh fitur antarmuka pengguna yang direncanakan telah berhasil
diimplementasikan secara penuh, mulai dari portal autentikasi glassmorphism yang
elegan, modul pengaduan lengkap dengan fitur GPS dan upload foto, peta interaktif
sebaran laporan, hingga dasbor administrasi multi-level.

Kode frontend yang dihasilkan ($\pm$2.900 baris JSX + CSS) bersifat
\textbf{production-ready} dengan design system yang konsisten, manajemen state
yang terstruktur, dan integrasi penuh dengan semua endpoint API backend yang telah
dibangun sebelumnya.

Proyek GovSmart kini memiliki fondasi frontend yang solid dan siap untuk
memasuki fase pengujian, optimasi, dan deployment pada minggu-minggu berikutnya.

\begin{infobox}[Catatan Tambahan]
  Seluruh kode sumber proyek GovSmart tersedia di direktori kerja
  \texttt{GovSmart-main/} dengan struktur monorepo yang memisahkan
  \texttt{frontend/} dan \texttt{backend/} secara jelas. Untuk menjalankan
  proyek, cukup menjalankan \texttt{docker-compose up} dari direktori root.
\end{infobox}

\vspace{1em}
\noindent\rule{\textwidth}{0.5pt}\\[0.3em]
{\small\color{gray}
Laporan ini dibuat sebagai bagian dari dokumentasi akademik\\
Workshop Pemrograman Framework --- Teknik Informatika}

\end{document}
